\documentclass[11pt,a4paper]{article}
% preamble.tex - shared house preamble for /explain-projects documents.
% Every package here is VERIFIED present in this machine's TeX Live install.
% Do not add packages without running: kpsewhich <pkg>.sty
% Notably ABSENT (do not \usepackage them, the build will fail):
%   algorithm2e, algorithm, algorithmic, algpseudocode  (texlive-science is not installed)
%   -> use the \begin{pseudocode} tcolorbox environment defined below instead.
%   lmodern, charter, mathptmx, times, mathpazo, courier, newtx*, tgtermes
%   -> their .sty files may resolve but the FONT FILES are absent, so the build
%      dies with "I can't find file `bchr8t'" or similar. Verified 2026-07-16 by
%      test-compiling each one. Only Computer Modern (the default) and helvet
%      actually render on this box. Do not "upgrade" the font stack without
%      test-compiling first: kpsewhich lying about a font is the exact trap here.
%
% To get a nicer font stack + algorithm2e, install:
%   sudo apt install texlive-fonts-recommended texlive-science
% Until then, this preamble deliberately uses the default font.
%
% Usage from a document:
%   \documentclass[11pt,a4paper]{article}
%   % preamble.tex - shared house preamble for /explain-projects documents.
% Every package here is VERIFIED present in this machine's TeX Live install.
% Do not add packages without running: kpsewhich <pkg>.sty
% Notably ABSENT (do not \usepackage them, the build will fail):
%   algorithm2e, algorithm, algorithmic, algpseudocode  (texlive-science is not installed)
%   -> use the \begin{pseudocode} tcolorbox environment defined below instead.
%   lmodern, charter, mathptmx, times, mathpazo, courier, newtx*, tgtermes
%   -> their .sty files may resolve but the FONT FILES are absent, so the build
%      dies with "I can't find file `bchr8t'" or similar. Verified 2026-07-16 by
%      test-compiling each one. Only Computer Modern (the default) and helvet
%      actually render on this box. Do not "upgrade" the font stack without
%      test-compiling first: kpsewhich lying about a font is the exact trap here.
%
% To get a nicer font stack + algorithm2e, install:
%   sudo apt install texlive-fonts-recommended texlive-science
% Until then, this preamble deliberately uses the default font.
%
% Usage from a document:
%   \documentclass[11pt,a4paper]{article}
%   % preamble.tex - shared house preamble for /explain-projects documents.
% Every package here is VERIFIED present in this machine's TeX Live install.
% Do not add packages without running: kpsewhich <pkg>.sty
% Notably ABSENT (do not \usepackage them, the build will fail):
%   algorithm2e, algorithm, algorithmic, algpseudocode  (texlive-science is not installed)
%   -> use the \begin{pseudocode} tcolorbox environment defined below instead.
%   lmodern, charter, mathptmx, times, mathpazo, courier, newtx*, tgtermes
%   -> their .sty files may resolve but the FONT FILES are absent, so the build
%      dies with "I can't find file `bchr8t'" or similar. Verified 2026-07-16 by
%      test-compiling each one. Only Computer Modern (the default) and helvet
%      actually render on this box. Do not "upgrade" the font stack without
%      test-compiling first: kpsewhich lying about a font is the exact trap here.
%
% To get a nicer font stack + algorithm2e, install:
%   sudo apt install texlive-fonts-recommended texlive-science
% Until then, this preamble deliberately uses the default font.
%
% Usage from a document:
%   \documentclass[11pt,a4paper]{article}
%   \input{preamble}          % this file is copied next to the .tex

\usepackage[T1]{fontenc}
\usepackage[utf8]{inputenc}
% Font: Computer Modern (default). See the note above before changing this.
\usepackage[a4paper,margin=2.4cm,headheight=14pt]{geometry}
% microtype WITHOUT font expansion. This is load-bearing, do not "restore" it:
% cm-super is not installed, so T1 Computer Modern resolves to BITMAP fonts, and
% pdfTeX's auto-expansion needs scalable ones. With expansion on, any document
% long enough to reach a second page dies with
%   "pdfTeX error (font expansion): auto expansion is only possible with scalable fonts"
% Protrusion (the bigger visual win) still works and stays on.
% `sudo apt install cm-super` would lift this; until then, leave expansion off.
\usepackage[protrusion=true,expansion=false]{microtype}
\usepackage{graphicx}
\usepackage{xcolor}
\usepackage{booktabs}
\usepackage{enumitem}
\usepackage{fancyhdr}
\usepackage{titlesec}
\usepackage{tcolorbox}
\usepackage{listings}
\usepackage{float}
\usepackage{caption}
\usepackage{tikz}
\usepackage{forest}
\usepackage{pgfplots}
\pgfplotsset{compat=1.18}
\usetikzlibrary{arrows.meta,positioning,shapes.geometric,fit,backgrounds,calc,decorations.pathreplacing}
\tcbuselibrary{skins,breakable}
\usepackage[hidelinks,bookmarks=true,bookmarksnumbered=true]{hyperref}

% ===== palette ===========================================================
\definecolor{inkblue}{HTML}{1F3A5F}
\definecolor{accent}{HTML}{2E7DAF}
\definecolor{softgrey}{HTML}{F4F6F8}
\definecolor{linegrey}{HTML}{C8D0D8}
\definecolor{codebg}{HTML}{F7F7F4}
\definecolor{warnred}{HTML}{A33B3B}
\definecolor{okgreen}{HTML}{2E7D53}

% ===== headings ==========================================================
\titleformat{\section}{\Large\bfseries\color{inkblue}}{\thesection}{0.7em}{}
\titleformat{\subsection}{\large\bfseries\color{inkblue}}{\thesubsection}{0.6em}{}
\titleformat{\subsubsection}{\normalsize\bfseries\color{accent}}{\thesubsubsection}{0.5em}{}

\pagestyle{fancy}
\fancyhf{}
\fancyhead[L]{\small\color{inkblue}\nouppercase{\leftmark}}
\fancyhead[R]{\small\color{inkblue}\thepage}
\renewcommand{\headrulewidth}{0.4pt}
\renewcommand{\headrule}{\hbox to\headwidth{\color{linegrey}\leaders\hrule height \headrulewidth\hfill}}

% ===== code ==============================================================
\lstdefinestyle{house}{
  backgroundcolor=\color{codebg},
  basicstyle=\ttfamily\footnotesize,
  keywordstyle=\color{inkblue}\bfseries,
  commentstyle=\color{gray}\itshape,
  stringstyle=\color{okgreen},
  numbers=left, numberstyle=\tiny\color{gray}, numbersep=8pt,
  frame=single, rulecolor=\color{linegrey},
  breaklines=true, breakatwhitespace=false,
  showstringspaces=false, tabsize=2,
  captionpos=b, xleftmargin=14pt, framexleftmargin=14pt,
  columns=fullflexible, keepspaces=true,
  upquote=true  % render ' and " as straight quotes, not curly
}
\lstset{style=house}

% ===== callout boxes =====================================================
% \begin{keyidea}{Title} ... \end{keyidea}   - a concept a beginner must grasp
\newtcolorbox{keyidea}[1]{
  enhanced, breakable, colback=softgrey, colframe=accent,
  boxrule=0.4pt, arc=2pt, left=8pt, right=8pt, top=6pt, bottom=6pt,
  title={\bfseries #1}, coltitle=white, colbacktitle=accent, fonttitle=\small
}

% \begin{jargon}{Term} definition \end{jargon}  - define a term at first use
\newtcolorbox{jargon}[1]{
  enhanced, breakable, colback=white, colframe=linegrey,
  boxrule=0.4pt, arc=2pt, left=8pt, right=8pt, top=6pt, bottom=6pt,
  title={\bfseries Jargon: #1}, coltitle=inkblue, colbacktitle=softgrey, fonttitle=\small
}

% \begin{honest}{Title} ... \end{honest}  - brutally honest finding (rule 18)
\newtcolorbox{honest}[1]{
  enhanced, breakable, colback=white, colframe=warnred,
  boxrule=0.6pt, arc=2pt, left=8pt, right=8pt, top=6pt, bottom=6pt,
  title={\bfseries #1}, coltitle=white, colbacktitle=warnred, fonttitle=\small
}

% \begin{provenance}{path/to/file} ... \end{provenance}
% Where a file came from: hand-written, generated by a tool, scaffolded by a
% framework, copied or adapted from an upstream project, ported from an older
% module. State the EVIDENCE (a header comment, a git log entry, an upstream
% URL, a generator command) and label the confidence. If no evidence exists,
% say so in the box. Never guess a lineage into existence.
\newtcolorbox{provenance}[1]{
  enhanced, breakable, colback=white, colframe=okgreen,
  boxrule=0.4pt, arc=2pt, left=8pt, right=8pt, top=6pt, bottom=6pt,
  title={\bfseries Where this came from: #1}, coltitle=white,
  colbacktitle=okgreen, fonttitle=\small
}

% \begin{pseudocode}{Caption} ... \end{pseudocode}
% Stands in for algorithm2e, which is NOT installed. Use \Step / \Indent inside.
\newtcolorbox{pseudocode}[1]{
  enhanced, breakable, colback=codebg, colframe=inkblue,
  boxrule=0.4pt, arc=2pt, left=10pt, right=8pt, top=6pt, bottom=6pt,
  fontupper=\ttfamily\small,
  title={\bfseries\rmfamily #1}, coltitle=white, colbacktitle=inkblue, fonttitle=\small
}
\newcommand{\Step}[1]{\par\noindent #1}
\newcommand{\Indent}[1]{\par\noindent\hspace*{2em}#1}

% ===== tikz house styles =================================================
\tikzset{
  box/.style   = {rectangle, rounded corners=2pt, draw=inkblue, fill=softgrey,
                  text=inkblue, align=center, inner sep=6pt, minimum height=9mm,
                  font=\small},
  extbox/.style= {rectangle, rounded corners=2pt, draw=linegrey, fill=white,
                  text=black!70, align=center, inner sep=6pt, minimum height=9mm,
                  font=\small, dashed},
  store/.style = {cylinder, shape border rotate=90, aspect=0.25, draw=accent,
                  fill=white, text=inkblue, align=center, inner sep=5pt,
                  font=\small},
  dec/.style   = {diamond, aspect=1.6, draw=accent, fill=white, text=inkblue,
                  align=center, inner sep=2pt, font=\scriptsize},
  flow/.style  = {-{Stealth[length=2.4mm]}, draw=inkblue, thick},
  dflow/.style = {-{Stealth[length=2.4mm]}, draw=accent, thick, dashed},
  % align=center is load-bearing: without it, a \\ inside an edge label throws
  % the misleading "perhaps a missing \item" error.
  lbl/.style   = {font=\scriptsize\itshape, text=black!65, fill=white,
                  inner sep=2pt, align=center}
}

% forest default for directory trees
\forestset{
  dirtree/.style={
    for tree={
      font=\ttfamily\small, grow'=0, child anchor=west, parent anchor=south,
      anchor=west, calign=first, inner sep=2pt,
      edge path={\noexpand\path [draw, \forestoption{edge}]
        (!u.south west) +(9pt,0) |- (.child anchor) \forestoption{edge label};},
      before typesetting nodes={if n=1{insert before={[,phantom]}}{}},
      fit=band, before computing xy={l=15pt}
    }
  }
}

% ===== misc ==============================================================
\captionsetup{font=small, labelfont={bf,color=inkblue}}
\setlist[itemize]{leftmargin=1.2em, itemsep=2pt, topsep=3pt}
\setlist[enumerate]{leftmargin=1.4em, itemsep=2pt, topsep=3pt}
\newcommand{\file}[1]{\texttt{\small #1}}
\newcommand{\code}[1]{\texttt{\small #1}}
\setcounter{tocdepth}{2}
\setcounter{secnumdepth}{3}
          % this file is copied next to the .tex

\usepackage[T1]{fontenc}
\usepackage[utf8]{inputenc}
% Font: Computer Modern (default). See the note above before changing this.
\usepackage[a4paper,margin=2.4cm,headheight=14pt]{geometry}
% microtype WITHOUT font expansion. This is load-bearing, do not "restore" it:
% cm-super is not installed, so T1 Computer Modern resolves to BITMAP fonts, and
% pdfTeX's auto-expansion needs scalable ones. With expansion on, any document
% long enough to reach a second page dies with
%   "pdfTeX error (font expansion): auto expansion is only possible with scalable fonts"
% Protrusion (the bigger visual win) still works and stays on.
% `sudo apt install cm-super` would lift this; until then, leave expansion off.
\usepackage[protrusion=true,expansion=false]{microtype}
\usepackage{graphicx}
\usepackage{xcolor}
\usepackage{booktabs}
\usepackage{enumitem}
\usepackage{fancyhdr}
\usepackage{titlesec}
\usepackage{tcolorbox}
\usepackage{listings}
\usepackage{float}
\usepackage{caption}
\usepackage{tikz}
\usepackage{forest}
\usepackage{pgfplots}
\pgfplotsset{compat=1.18}
\usetikzlibrary{arrows.meta,positioning,shapes.geometric,fit,backgrounds,calc,decorations.pathreplacing}
\tcbuselibrary{skins,breakable}
\usepackage[hidelinks,bookmarks=true,bookmarksnumbered=true]{hyperref}

% ===== palette ===========================================================
\definecolor{inkblue}{HTML}{1F3A5F}
\definecolor{accent}{HTML}{2E7DAF}
\definecolor{softgrey}{HTML}{F4F6F8}
\definecolor{linegrey}{HTML}{C8D0D8}
\definecolor{codebg}{HTML}{F7F7F4}
\definecolor{warnred}{HTML}{A33B3B}
\definecolor{okgreen}{HTML}{2E7D53}

% ===== headings ==========================================================
\titleformat{\section}{\Large\bfseries\color{inkblue}}{\thesection}{0.7em}{}
\titleformat{\subsection}{\large\bfseries\color{inkblue}}{\thesubsection}{0.6em}{}
\titleformat{\subsubsection}{\normalsize\bfseries\color{accent}}{\thesubsubsection}{0.5em}{}

\pagestyle{fancy}
\fancyhf{}
\fancyhead[L]{\small\color{inkblue}\nouppercase{\leftmark}}
\fancyhead[R]{\small\color{inkblue}\thepage}
\renewcommand{\headrulewidth}{0.4pt}
\renewcommand{\headrule}{\hbox to\headwidth{\color{linegrey}\leaders\hrule height \headrulewidth\hfill}}

% ===== code ==============================================================
\lstdefinestyle{house}{
  backgroundcolor=\color{codebg},
  basicstyle=\ttfamily\footnotesize,
  keywordstyle=\color{inkblue}\bfseries,
  commentstyle=\color{gray}\itshape,
  stringstyle=\color{okgreen},
  numbers=left, numberstyle=\tiny\color{gray}, numbersep=8pt,
  frame=single, rulecolor=\color{linegrey},
  breaklines=true, breakatwhitespace=false,
  showstringspaces=false, tabsize=2,
  captionpos=b, xleftmargin=14pt, framexleftmargin=14pt,
  columns=fullflexible, keepspaces=true,
  upquote=true  % render ' and " as straight quotes, not curly
}
\lstset{style=house}

% ===== callout boxes =====================================================
% \begin{keyidea}{Title} ... \end{keyidea}   - a concept a beginner must grasp
\newtcolorbox{keyidea}[1]{
  enhanced, breakable, colback=softgrey, colframe=accent,
  boxrule=0.4pt, arc=2pt, left=8pt, right=8pt, top=6pt, bottom=6pt,
  title={\bfseries #1}, coltitle=white, colbacktitle=accent, fonttitle=\small
}

% \begin{jargon}{Term} definition \end{jargon}  - define a term at first use
\newtcolorbox{jargon}[1]{
  enhanced, breakable, colback=white, colframe=linegrey,
  boxrule=0.4pt, arc=2pt, left=8pt, right=8pt, top=6pt, bottom=6pt,
  title={\bfseries Jargon: #1}, coltitle=inkblue, colbacktitle=softgrey, fonttitle=\small
}

% \begin{honest}{Title} ... \end{honest}  - brutally honest finding (rule 18)
\newtcolorbox{honest}[1]{
  enhanced, breakable, colback=white, colframe=warnred,
  boxrule=0.6pt, arc=2pt, left=8pt, right=8pt, top=6pt, bottom=6pt,
  title={\bfseries #1}, coltitle=white, colbacktitle=warnred, fonttitle=\small
}

% \begin{provenance}{path/to/file} ... \end{provenance}
% Where a file came from: hand-written, generated by a tool, scaffolded by a
% framework, copied or adapted from an upstream project, ported from an older
% module. State the EVIDENCE (a header comment, a git log entry, an upstream
% URL, a generator command) and label the confidence. If no evidence exists,
% say so in the box. Never guess a lineage into existence.
\newtcolorbox{provenance}[1]{
  enhanced, breakable, colback=white, colframe=okgreen,
  boxrule=0.4pt, arc=2pt, left=8pt, right=8pt, top=6pt, bottom=6pt,
  title={\bfseries Where this came from: #1}, coltitle=white,
  colbacktitle=okgreen, fonttitle=\small
}

% \begin{pseudocode}{Caption} ... \end{pseudocode}
% Stands in for algorithm2e, which is NOT installed. Use \Step / \Indent inside.
\newtcolorbox{pseudocode}[1]{
  enhanced, breakable, colback=codebg, colframe=inkblue,
  boxrule=0.4pt, arc=2pt, left=10pt, right=8pt, top=6pt, bottom=6pt,
  fontupper=\ttfamily\small,
  title={\bfseries\rmfamily #1}, coltitle=white, colbacktitle=inkblue, fonttitle=\small
}
\newcommand{\Step}[1]{\par\noindent #1}
\newcommand{\Indent}[1]{\par\noindent\hspace*{2em}#1}

% ===== tikz house styles =================================================
\tikzset{
  box/.style   = {rectangle, rounded corners=2pt, draw=inkblue, fill=softgrey,
                  text=inkblue, align=center, inner sep=6pt, minimum height=9mm,
                  font=\small},
  extbox/.style= {rectangle, rounded corners=2pt, draw=linegrey, fill=white,
                  text=black!70, align=center, inner sep=6pt, minimum height=9mm,
                  font=\small, dashed},
  store/.style = {cylinder, shape border rotate=90, aspect=0.25, draw=accent,
                  fill=white, text=inkblue, align=center, inner sep=5pt,
                  font=\small},
  dec/.style   = {diamond, aspect=1.6, draw=accent, fill=white, text=inkblue,
                  align=center, inner sep=2pt, font=\scriptsize},
  flow/.style  = {-{Stealth[length=2.4mm]}, draw=inkblue, thick},
  dflow/.style = {-{Stealth[length=2.4mm]}, draw=accent, thick, dashed},
  % align=center is load-bearing: without it, a \\ inside an edge label throws
  % the misleading "perhaps a missing \item" error.
  lbl/.style   = {font=\scriptsize\itshape, text=black!65, fill=white,
                  inner sep=2pt, align=center}
}

% forest default for directory trees
\forestset{
  dirtree/.style={
    for tree={
      font=\ttfamily\small, grow'=0, child anchor=west, parent anchor=south,
      anchor=west, calign=first, inner sep=2pt,
      edge path={\noexpand\path [draw, \forestoption{edge}]
        (!u.south west) +(9pt,0) |- (.child anchor) \forestoption{edge label};},
      before typesetting nodes={if n=1{insert before={[,phantom]}}{}},
      fit=band, before computing xy={l=15pt}
    }
  }
}

% ===== misc ==============================================================
\captionsetup{font=small, labelfont={bf,color=inkblue}}
\setlist[itemize]{leftmargin=1.2em, itemsep=2pt, topsep=3pt}
\setlist[enumerate]{leftmargin=1.4em, itemsep=2pt, topsep=3pt}
\newcommand{\file}[1]{\texttt{\small #1}}
\newcommand{\code}[1]{\texttt{\small #1}}
\setcounter{tocdepth}{2}
\setcounter{secnumdepth}{3}
          % this file is copied next to the .tex

\usepackage[T1]{fontenc}
\usepackage[utf8]{inputenc}
% Font: Computer Modern (default). See the note above before changing this.
\usepackage[a4paper,margin=2.4cm,headheight=14pt]{geometry}
% microtype WITHOUT font expansion. This is load-bearing, do not "restore" it:
% cm-super is not installed, so T1 Computer Modern resolves to BITMAP fonts, and
% pdfTeX's auto-expansion needs scalable ones. With expansion on, any document
% long enough to reach a second page dies with
%   "pdfTeX error (font expansion): auto expansion is only possible with scalable fonts"
% Protrusion (the bigger visual win) still works and stays on.
% `sudo apt install cm-super` would lift this; until then, leave expansion off.
\usepackage[protrusion=true,expansion=false]{microtype}
\usepackage{graphicx}
\usepackage{xcolor}
\usepackage{booktabs}
\usepackage{enumitem}
\usepackage{fancyhdr}
\usepackage{titlesec}
\usepackage{tcolorbox}
\usepackage{listings}
\usepackage{float}
\usepackage{caption}
\usepackage{tikz}
\usepackage{forest}
\usepackage{pgfplots}
\pgfplotsset{compat=1.18}
\usetikzlibrary{arrows.meta,positioning,shapes.geometric,fit,backgrounds,calc,decorations.pathreplacing}
\tcbuselibrary{skins,breakable}
\usepackage[hidelinks,bookmarks=true,bookmarksnumbered=true]{hyperref}

% ===== palette ===========================================================
\definecolor{inkblue}{HTML}{1F3A5F}
\definecolor{accent}{HTML}{2E7DAF}
\definecolor{softgrey}{HTML}{F4F6F8}
\definecolor{linegrey}{HTML}{C8D0D8}
\definecolor{codebg}{HTML}{F7F7F4}
\definecolor{warnred}{HTML}{A33B3B}
\definecolor{okgreen}{HTML}{2E7D53}

% ===== headings ==========================================================
\titleformat{\section}{\Large\bfseries\color{inkblue}}{\thesection}{0.7em}{}
\titleformat{\subsection}{\large\bfseries\color{inkblue}}{\thesubsection}{0.6em}{}
\titleformat{\subsubsection}{\normalsize\bfseries\color{accent}}{\thesubsubsection}{0.5em}{}

\pagestyle{fancy}
\fancyhf{}
\fancyhead[L]{\small\color{inkblue}\nouppercase{\leftmark}}
\fancyhead[R]{\small\color{inkblue}\thepage}
\renewcommand{\headrulewidth}{0.4pt}
\renewcommand{\headrule}{\hbox to\headwidth{\color{linegrey}\leaders\hrule height \headrulewidth\hfill}}

% ===== code ==============================================================
\lstdefinestyle{house}{
  backgroundcolor=\color{codebg},
  basicstyle=\ttfamily\footnotesize,
  keywordstyle=\color{inkblue}\bfseries,
  commentstyle=\color{gray}\itshape,
  stringstyle=\color{okgreen},
  numbers=left, numberstyle=\tiny\color{gray}, numbersep=8pt,
  frame=single, rulecolor=\color{linegrey},
  breaklines=true, breakatwhitespace=false,
  showstringspaces=false, tabsize=2,
  captionpos=b, xleftmargin=14pt, framexleftmargin=14pt,
  columns=fullflexible, keepspaces=true,
  upquote=true  % render ' and " as straight quotes, not curly
}
\lstset{style=house}

% ===== callout boxes =====================================================
% \begin{keyidea}{Title} ... \end{keyidea}   - a concept a beginner must grasp
\newtcolorbox{keyidea}[1]{
  enhanced, breakable, colback=softgrey, colframe=accent,
  boxrule=0.4pt, arc=2pt, left=8pt, right=8pt, top=6pt, bottom=6pt,
  title={\bfseries #1}, coltitle=white, colbacktitle=accent, fonttitle=\small
}

% \begin{jargon}{Term} definition \end{jargon}  - define a term at first use
\newtcolorbox{jargon}[1]{
  enhanced, breakable, colback=white, colframe=linegrey,
  boxrule=0.4pt, arc=2pt, left=8pt, right=8pt, top=6pt, bottom=6pt,
  title={\bfseries Jargon: #1}, coltitle=inkblue, colbacktitle=softgrey, fonttitle=\small
}

% \begin{honest}{Title} ... \end{honest}  - brutally honest finding (rule 18)
\newtcolorbox{honest}[1]{
  enhanced, breakable, colback=white, colframe=warnred,
  boxrule=0.6pt, arc=2pt, left=8pt, right=8pt, top=6pt, bottom=6pt,
  title={\bfseries #1}, coltitle=white, colbacktitle=warnred, fonttitle=\small
}

% \begin{provenance}{path/to/file} ... \end{provenance}
% Where a file came from: hand-written, generated by a tool, scaffolded by a
% framework, copied or adapted from an upstream project, ported from an older
% module. State the EVIDENCE (a header comment, a git log entry, an upstream
% URL, a generator command) and label the confidence. If no evidence exists,
% say so in the box. Never guess a lineage into existence.
\newtcolorbox{provenance}[1]{
  enhanced, breakable, colback=white, colframe=okgreen,
  boxrule=0.4pt, arc=2pt, left=8pt, right=8pt, top=6pt, bottom=6pt,
  title={\bfseries Where this came from: #1}, coltitle=white,
  colbacktitle=okgreen, fonttitle=\small
}

% \begin{pseudocode}{Caption} ... \end{pseudocode}
% Stands in for algorithm2e, which is NOT installed. Use \Step / \Indent inside.
\newtcolorbox{pseudocode}[1]{
  enhanced, breakable, colback=codebg, colframe=inkblue,
  boxrule=0.4pt, arc=2pt, left=10pt, right=8pt, top=6pt, bottom=6pt,
  fontupper=\ttfamily\small,
  title={\bfseries\rmfamily #1}, coltitle=white, colbacktitle=inkblue, fonttitle=\small
}
\newcommand{\Step}[1]{\par\noindent #1}
\newcommand{\Indent}[1]{\par\noindent\hspace*{2em}#1}

% ===== tikz house styles =================================================
\tikzset{
  box/.style   = {rectangle, rounded corners=2pt, draw=inkblue, fill=softgrey,
                  text=inkblue, align=center, inner sep=6pt, minimum height=9mm,
                  font=\small},
  extbox/.style= {rectangle, rounded corners=2pt, draw=linegrey, fill=white,
                  text=black!70, align=center, inner sep=6pt, minimum height=9mm,
                  font=\small, dashed},
  store/.style = {cylinder, shape border rotate=90, aspect=0.25, draw=accent,
                  fill=white, text=inkblue, align=center, inner sep=5pt,
                  font=\small},
  dec/.style   = {diamond, aspect=1.6, draw=accent, fill=white, text=inkblue,
                  align=center, inner sep=2pt, font=\scriptsize},
  flow/.style  = {-{Stealth[length=2.4mm]}, draw=inkblue, thick},
  dflow/.style = {-{Stealth[length=2.4mm]}, draw=accent, thick, dashed},
  % align=center is load-bearing: without it, a \\ inside an edge label throws
  % the misleading "perhaps a missing \item" error.
  lbl/.style   = {font=\scriptsize\itshape, text=black!65, fill=white,
                  inner sep=2pt, align=center}
}

% forest default for directory trees
\forestset{
  dirtree/.style={
    for tree={
      font=\ttfamily\small, grow'=0, child anchor=west, parent anchor=south,
      anchor=west, calign=first, inner sep=2pt,
      edge path={\noexpand\path [draw, \forestoption{edge}]
        (!u.south west) +(9pt,0) |- (.child anchor) \forestoption{edge label};},
      before typesetting nodes={if n=1{insert before={[,phantom]}}{}},
      fit=band, before computing xy={l=15pt}
    }
  }
}

% ===== misc ==============================================================
\captionsetup{font=small, labelfont={bf,color=inkblue}}
\setlist[itemize]{leftmargin=1.2em, itemsep=2pt, topsep=3pt}
\setlist[enumerate]{leftmargin=1.4em, itemsep=2pt, topsep=3pt}
\newcommand{\file}[1]{\texttt{\small #1}}
\newcommand{\code}[1]{\texttt{\small #1}}
\setcounter{tocdepth}{2}
\setcounter{secnumdepth}{3}


\DeclareUnicodeCharacter{2500}{-}
\DeclareUnicodeCharacter{2550}{=}
\lstset{extendedchars=true,
  literate={·}{{\textperiodcentered}}1
           {─}{{-}}1
           {═}{{=}}1
}

\title{Salman Adnan Site\\[2mm]\large Brief: the whole project in one sitting}
\author{Digitalise Agency}
\date{2026-08-11}

\begin{document}
\maketitle
\tableofcontents
\newpage

% ===========================================================================
\section{What it is, and two things to know first}
% ===========================================================================

This repository holds a personal brand website for \textbf{salmanadnan.com}: one
page introducing Salman Adnan, showing six client projects, and asking the visitor
to book a thirty-minute call.

Six files, 1{,}180 lines, no framework, no build step, no package manager, no
database, no server code. That is unusual enough in 2026 to be worth stating as a
decision rather than an absence.

\begin{keyidea}{First: this is not the page being served at that address}
The design here was finished, then deliberately parked. On 2026-07-29 one commit
moved every file into a folder called \file{new-design-reference/} and renamed the
publishing workflow so it could no longer run. Its message says why: the live site
is served from a different repository.

Fetching \code{https://salmanadnan.com/} on 2026-08-11 confirms it. The address
answers with \code{200} from the same VPS this project deploys to, and serves a
page with a different title, a proper social preview image, and none of this
version's distinctive markup. Read this repository as a design reference, never as
the live site.
\end{keyidea}

\begin{keyidea}{Second: the design was never finished, and it says so on screen}
The left half of the hero, which is half of everything visible when the page loads,
contains no photograph. It contains a giant letter \code{S} on an orange gradient
and the words \emph{``Drop Salman's photo here / assets/salman.jpg''}.

That is not a code comment. It is content, deliberately styled and positioned, and
any visitor would read it. On a personal brand page, whose entire proposition is a
person, the missing element is the person. This is why the honest description of
this project is ``unfinished'' rather than ``finished and parked''.
\end{keyidea}

\begin{figure}[H]\centering
\begin{tikzpicture}[node distance=10mm and 18mm]
  \node[box]                (v)  {A visitor\\\scriptsize from a link or a referral};
  \node[box, right=of v]    (p)  {The page\\\scriptsize \file{index.html}};
  \node[extbox, right=of p] (cal){cal.com\\\scriptsize \code{salman-adnan/meeting}};
  \node[box, below=of p]    (out){A booked call\\\scriptsize the only outcome that counts};
  \draw[flow]  (v) -- node[lbl]{opens} (p);
  \draw[dflow] (p) -- node[lbl]{Book a call} (cal);
  \draw[flow]  (cal) |- node[lbl, pos=0.25]{picks a time} (out);
\end{tikzpicture}
\caption{The entire purpose. Dashed boxes are services the agency does not run.}
\end{figure}

\subsection{Read it beside its twin}

There is a near-identical sibling in this workspace, \code{ai-digitalise}, built the
same week from the same skeleton: the same six file paths, the same three-file
architecture, the same publishing workflow and the same fate. They are not copies;
every substantive file differs. This one is a personal page in orange, that one an
agency service page in blue.

Comparing them teaches more than reading either alone, and the Deep Dive points out
each instructive difference where it arises.

% ===========================================================================
\section{The mental model}
% ===========================================================================

\begin{figure}[H]\centering
\begin{tikzpicture}[node distance=11mm and 14mm]
  \node[box]                 (html) {\file{index.html}\\\scriptsize 395 lines: words and structure};
  \node[box, right=of html]  (css)  {\file{css/site.css}\\\scriptsize 575 lines: what it looks like};
  \node[box, right=of css]   (js)   {\file{js/main.js}\\\scriptsize 64 lines: what moves};
  \node[extbox, below=of css](br)   {The visitor's browser\\\scriptsize assembles all three};
  \draw[flow] (html) -- node[lbl]{names} (css);
  \draw[flow] (css)  -- node[lbl]{then} (js);
  \draw[flow] (html) |- (br);
  \draw[flow] (css)  -- (br);
  \draw[flow] (js)   |- (br);
\end{tikzpicture}
\caption{Nothing runs on a server. Everything happens on the visitor's machine.}
\end{figure}

\begin{keyidea}{The arrows only point one way, and one of them is already broken}
\file{index.html} names the other two files. Neither knows anything about it: the
stylesheet styles class names it hopes exist, and the JavaScript looks for elements
that may not be there, guarding every lookup so it does nothing quietly when they
are missing.

Nothing checks that the three agree, and in one place they do not. Three stylesheet
rules target a class the HTML never uses, so they have never once applied. The page
looks correct anyway, because the HTML achieves the same layout by a different
means. Nothing anywhere reported it. That is what ``silent failure'' means in
practice.
\end{keyidea}

% ===========================================================================
\section{How it works, end to end}
% ===========================================================================

\begin{figure}[H]\centering
\begin{tikzpicture}[node distance=8mm and 12mm]
  \node[box]             (a) {Browser downloads\\\file{index.html}};
  \node[box, right=of a] (b) {Requests fonts\\and the stylesheet};
  \node[box, right=of b] (c) {Draws the page\\top to bottom};
  \node[box, below=of c] (d) {Reaches \code{<script>}\\at the very end};
  \node[box, left=of d]  (e) {\file{main.js} runs,\\attaches watchers};
  \node[box, left=of e]  (f) {Visitor scrolls,\\watchers fire};
  \draw[flow] (a) -- (b); \draw[flow] (b) -- (c); \draw[flow] (c) -- (d);
  \draw[flow] (d) -- (e); \draw[flow] (e) -- (f);
\end{tikzpicture}
\caption{One page view, start to finish.}
\end{figure}

The page stacks six bands: a fixed navigation bar; a two-column hero with the photo
frame on the left, three drifting frosted badges over it, and the name, tagline,
buttons and three small statistics on the right; a grid of six client projects;
four large counting numbers; two case-study cards; a booking panel; and a footer.

Three things react to scrolling, and that is the whole of the interactivity: the
navigation bar gains a frosted background past 40 pixels of scroll, the four large
numbers count up from zero over 2.2 seconds when they come into view, and most
blocks fade upward as they arrive.

There is no form, nothing is saved, and nothing is sent anywhere. The only outbound
action is following the booking link, which opens a cal.com calendar in a new tab.

\begin{honest}{Nothing measures whether the page works}
No analytics of any kind: no counter, no tag manager, no event when somebody presses
``Book a call''. The one question the page exists to answer cannot be answered from
the page. The live salmanadnan.com does carry a full social-preview setup this
version lacks, so whoever built the replacement thought about distribution in a way
this design did not.
\end{honest}

% ===========================================================================
\section{The main components}
% ===========================================================================

\begin{table}[H]\centering\small
\begin{tabular}{@{}llrp{5.4cm}@{}}
\toprule
File & Kind & Lines & What it does \\
\midrule
\file{new-design-reference/index.html} & HTML & 395 & Every word on the page and its structure, with all icons as inline drawings \\
\file{new-design-reference/css/site.css} & CSS & 575 & All appearance, layout, movement, and the two mobile layouts \\
\file{new-design-reference/js/main.js} & JavaScript & 64 & The three scroll behaviours, no dependencies \\
\file{new-design-reference/netlify.toml} & config & 11 & Hosting settings for a host this project does not use \\
\file{new-design-reference/assets/ASSETS.md} & notes & 32 & A checklist of eleven assets, none of which were produced \\
\file{.github/workflows-disabled/deploy.yml.disabled} & workflow & 103 & The publishing pipeline, switched off \\
\bottomrule
\end{tabular}
\caption{Every file git tracks. The Deep Dive walks all six line by line, in
sections 7.1 to 7.6.}
\end{table}

\subsection{\file{index.html}: the page}

The content and structure. Unusually self-contained: every icon is an SVG drawing
written directly into the file rather than loaded as an image, so the page needs no
image downloads at all.

Two habits are worth knowing before editing it. Appearance is sometimes written
directly onto elements with \code{style} attributes rather than left to the
stylesheet, and an inline style always wins. And the four large statistics keep
their real values in \code{data-} attributes with a visible \code{0}, so they read
as zero until JavaScript runs. Deep Dive section 7.1.

\subsection{\file{site.css}: the design system}

The most valuable file here, and the one worth borrowing. It opens with nineteen
named values, the ``tokens'', covering surfaces, text, accents, borders, typefaces,
radii, a motion curve and a content width. Its own second line records the lineage:
the tokens are shared with ai.digitalise.agency.

\begin{keyidea}{Why the token names are numbers rather than colours}
Calling a colour \code{-{}-s1} instead of \code{-{}-dark-navy} names what it is
\emph{for} (the surface one step up from the page) rather than how it looks. This
pair of repositories proves the point: the two stylesheets share every token
\emph{name} and differ in nearly every token \emph{value}, so one vocabulary
produced a blue agency site and an orange personal one with no renaming at all.
\end{keyidea}

It ends with the best thing in the project: a block honouring the visitor's
reduced-motion preference, which crushes every animation, forces the fade-in
elements to full opacity so they still appear, and explicitly stops the drifting
badges. That is more careful than the sibling's equivalent. Deep Dive section 7.2.

\subsection{\file{main.js}: the three behaviours}

Sixty-four lines, no dependencies, better written than its size suggests: a wrapper
keeping its names private, a guard at every lookup, a small shared helper built on
\code{IntersectionObserver} that three one-line calls reuse, and correct use of the
flags that keep scrolling smooth and stop each animation firing twice.

Its interface with the rest of the project is six strings. Break any one and the
page keeps working, minus one behaviour, silently. Deep Dive section 7.3.

\subsection{The three files about hosting}

\file{netlify.toml} declares four sensible security headers for a hosting service
this project does not use, and is byte for byte identical to the sibling's copy.
\file{ASSETS.md} lists the eleven assets needed to finish the design, none of which
exist. The workflow publishes to the agency's server when a commit subject begins
with \code{[deploy]}, and is switched off. Deep Dive sections 7.4 to 7.6.

% ===========================================================================
\section{How it is operated}
% ===========================================================================

\subsection{Looking at it locally}

\begin{lstlisting}[language=bash,caption={Serve it, do not open the file directly}]
cd new-design-reference
python3 -m http.server 8000
# then open http://localhost:8000/
\end{lstlisting}

Opening \file{index.html} straight off the disk renders it completely unstyled,
which looks like catastrophic breakage and is not. The page asks for its stylesheet
at \code{/css/site.css}, and the leading slash means ``from the top of the site'',
which off a disk is the root of your hard drive. This one character is the most
likely thing to make a newcomer think the project is broken.

\subsection{Publishing}

\begin{figure}[H]\centering
\begin{tikzpicture}[node distance=9mm and 13mm]
  \node[box]                (push) {Push to \code{main}};
  \node[dec, right=of push] (chk)  {subject starts\\\code{[deploy]}?};
  \node[box, right=of chk]  (ssh)  {Install key,\\fetch host key};
  \node[box, below=of ssh]  (sync) {\code{rsync} to the VPS};
  \node[box, left=of sync]  (ver)  {Check the site\\still answers};
  \node[box, left=of ver]   (stop) {Do nothing};
  \draw[flow] (push) -- (chk);
  \draw[flow] (chk) -- node[lbl]{yes} (ssh);
  \draw[flow] (chk) |- node[lbl,pos=0.3]{no} (stop);
  \draw[flow] (ssh) -- (sync);
  \draw[flow] (sync) -- (ver);
\end{tikzpicture}
\caption{The publishing pipeline as written. Switched off, twice over, since
2026-07-29.}
\end{figure}

The \code{[deploy]} marker must \emph{begin} the commit subject line. Writing ``fix
the header and [deploy]'' does nothing at all. This convention is shared by every
repository in this workspace.

The workflow is disabled two ways: its folder is \file{.github/workflows-disabled}
rather than \file{.github/workflows}, and the file ends \file{.disabled} rather than
\file{.yml}. Either alone would suffice, which tells you the intent was not to pause
it briefly.

\begin{honest}{Restoring it unchanged would take the live site down}
The \code{rsync} command copies the repository root, which since the move contains a
folder rather than a home page, and it copies with \code{-{}-delete}. So it would
replace the live salmanadnan.com with a directory listing's worth of nothing. The
correct source after the move is \code{new-design-reference/}.
\end{honest}

\subsection{Claude Code tooling}

\textbf{This project ships none}, which was checked rather than assumed: no
\file{.claude/} folder, no \file{CLAUDE.md}, no skill and no hook.

There is no README either, and that matters more than it sounds. The central fact
about this repository, that it is parked and the live site is elsewhere, exists only
in a commit message, invisible unless somebody thinks to run \code{git log}. The
same gap exists in the sibling.

The two PDFs in \file{docs/explainer/} are produced by a workspace-level command,
\code{/explain-projects}, living outside this repository.

% ===========================================================================
\section{Configuration and deployment at a glance}
% ===========================================================================

Nothing is configurable at run time: no settings file, no environment file. Three
values behave like configuration and are written into the source:

\begin{table}[H]\centering\small
\begin{tabular}{@{}lll@{}}
\toprule
Value & Where & Controls \\
\midrule
\code{https://cal.com/salman-adnan/meeting} & \file{index.html} line 362 & which calendar opens \\
\code{salman@digitalise.agency} & \file{index.html} line 386 & the footer contact \\
The nineteen tokens & \file{site.css} lines 5--25 & the entire colour scheme \\
\bottomrule
\end{tabular}
\caption{Values that behave like settings.}
\end{table}

The repository holds no passwords, keys or tokens, and should not: a page running
entirely in a visitor's browser can keep no secret. Three secrets live in GitHub and
are referenced by the workflow (\code{VPS\_HOST}, \code{VPS\_USER},
\code{VPS\_SSH\_KEY}); the server's IP address is also written into a comment in
that same file, which defeats the point of storing it as a secret.

The page has two run-time dependencies, both reached by the visitor's browser:
Google Fonts for two typefaces, and cal.com for the booking calendar.

% ===========================================================================
\section{Honest findings}
% ===========================================================================

\subsection{Unfinished}

\begin{itemize}
\item \textbf{The hero photo is a visible placeholder} reading ``Drop Salman's photo
  here''. The finding that matters most.
\item \textbf{None of the eleven assets in \file{ASSETS.md} exist}: no photograph,
  no project thumbnails, no case-study screenshots, no favicon, no social preview
  image. The live replacement site has at least the last of these.
\item \textbf{Seven stale comments} remain in \file{index.html}, some describing
  work that was done and some work that was not.
\end{itemize}

\subsection{The rebrand to orange was left half done}

\begin{honest}{Five rules still use the agency blue}
Commit \code{05210af} is titled ``Differentiate to orange brand''. The tokens were
duly changed. Five rules further down still hard-code the agency blue rather than
reading the token:

\begin{table}[H]\centering\small
\begin{tabular}{@{}llp{5.2cm}@{}}
\toprule
Line & Rule & What a visitor sees \\
\midrule
65  & \code{.btn-primary:hover} & an orange button casting a blue shadow \\
255--256 & \code{.hero-tag} & a blue-tinted pill containing orange text \\
337 & \code{.work-card:hover} & a project card outlined in blue on hover \\
435 & \code{.case-card:hover} & the same on the case-study cards \\
498 & \code{.book-wrap::before} & a blue glow above the booking panel \\
\bottomrule
\end{tabular}
\end{table}

Orange and blue are close to complementary, so the shadow under the button reads as
a mistake rather than a choice. All five are one-line fixes. Together they are a
clean illustration of what the token system is for: the rebrand changed the tokens
and could not reach the literals. Notably it \emph{did} reach the six thumbnail
gradients, which tells you it was done by eye rather than by search.
\end{honest}

\subsection{Dead code}

Three stylesheet rules for a reversal class the HTML never uses, which have never
applied and whose job is done by inline \code{order} properties instead; four tokens
defined and never read, including the agency blue that five other rules write out as
a literal; a class added by JavaScript that nothing styles; and a whole hosting
configuration file for a host that is not used.

\subsection{Content and consistency}

The three hero statistics duplicate three of the four in the metrics band, written
as plain text in one place and as data attributes in the other. \file{ASSETS.md}
knows this and says to update both, so it is a documented trade rather than an
accident, and it accidentally rescues the page from claiming zero of everything when
JavaScript does not run.

Two of the three instructions in \file{ASSETS.md} are stale. The six project cards
show a play button on hover and play nothing. The navigation logo is not a link.

\subsection{Accessibility}

\begin{itemize}
\item \textbf{The counting numbers ignore the reduced-motion preference.} The
  stylesheet honours it more carefully than the sibling does; the JavaScript never
  checks. Four lines fix it, given in Deep Dive section 7.3.
\item \textbf{Neither button has a focus state}, so keyboard navigation relies on
  whatever outline the browser draws, often nearly invisible on a dark background.
  These are the page's only interactive elements.
\item \textbf{The small orange uppercase labels are worth a contrast measurement.}
  This document does not have one and raises it rather than asserting a failure.
\end{itemize}

\subsection{One finding that belongs to a third repository}

\begin{honest}{The live ai.digitalise.agency claims this domain as its canonical}
The page live at ai.digitalise.agency declares
\code{<link rel="canonical" href="https://salmanadnan.com/">}, which asks search
engines to index this domain instead of that one. The live salmanadnan.com declares
its own canonical correctly, and both are served from the same VPS.

So the agency page was produced from this personal site's build and its canonical
link was never changed. Neither repository in this pair can fix it, because neither
contains the live build. It is recorded in both documents so whoever maintains the
portfolio repository finds it. Verified by fetching both addresses on 2026-08-11.
\end{honest}

% ===========================================================================
\section{The story, in brief}
% ===========================================================================

\subsection{What it delivered}

A complete, self-contained one-page personal brand website with no dependencies to
install; a hosting configuration setting four security headers; a GitHub Actions
workflow publishing over SSH with two retry loops and a post-deploy check; and an
asset checklist. What it did not deliver is a live site, or a photograph.

\subsection{When it was built}

Ten commits across two calendar days, 2026-07-28 and 2026-07-29, and nothing since.
The shape of that history is the interesting part: four commits on day one building
the thing, five consecutive commits on day two fighting the deploy, and one commit
switching it off.

This is calendar time, not effort. Git records when commits happened and nothing
about hours spent, and no document here records effort, so none is claimed.

\subsection{The four challenges that mattered}

Each is claimed because the repository carries a scar for it.

\begin{enumerate}
\item \textbf{SSH access to the server}, across three commits: secrets missing, then
  the wrong key, then a \code{deploy} user without permission to traverse into
  \code{/root}. That last is genuinely unobvious and its workaround survives in
  three \code{rsync} flags.
\item \textbf{GitHub's runners being refused}, across two more. One is titled
  ``Retry, runner IP lottery'', which is the diagnosis stated outright. This
  repository is where the evidence lives; the sibling has the same retry loops and
  no commit explaining them.
\item \textbf{Deciding what the brand was.} One commit exists solely to
  differentiate this site from the agency's blue, and it reached the tokens while
  missing five literals: the trace of a decision made after the design was written
  rather than before.
\item \textbf{Being overtaken.} The final commit is a retreat, not a fix. The page
  was finished and working when something else became the live site. The largest
  event here, and not technical.
\end{enumerate}

\subsection{Where it stands}

Built and unfinished: it renders, it is responsive, and the most prominent element
on the page is still a request for a photograph. Parked deliberately, with the
workflow disabled twice over. Superseded by a build that has a proper title, a
correct canonical link and a social preview image. One thing survived: the booking
link set here is the same one the live site uses today. Dormant since 2026-07-29.

\subsection{Do and do not}

\emph{From here on this is advice and judgement, not a record of decisions.}

\textbf{Do:}
\begin{itemize}
\item Read \code{git log} before touching anything: it holds the only statement of
  this repository's status.
\item Serve the folder over HTTP to view it, never open the file directly.
\item Read it beside \code{ai-digitalise}. The differences between the two are where
  the interesting decisions are.
\item Keep the \code{[deploy]} convention and both retry loops if the workflow is
  ever restored. They were paid for over five commits.
\item Treat the hero as the reusable asset: a two-column layout with fading edges
  and drifting frosted badges, in about eighty lines of CSS.
\end{itemize}

\textbf{Do not:}
\begin{itemize}
\item Re-enable the workflow without changing the \code{rsync} source path. It would
  copy a folder rather than a page, with \code{-{}-delete} set, and take the live site
  down.
\item Publish this design without the photograph.
\item Trust \file{ASSETS.md}. Two of its three instructions refer to things no
  longer in the code, or never there.
\item Fix the five blue literals by changing the token. The token is correct and the
  literals are wrong; changing \code{-{}-accent} back to blue would make five rules
  consistent and the entire site wrong.
\item Assume this repository is what visitors see. It is not.
\end{itemize}

\subsection{What would change, starting over}

\emph{Opinion, argued from the code.} Keep the no-build-step decision, the token
block, the three-behaviour JavaScript, the reduced-motion handling and the
two-column hero. For a page this size every one is right, and a framework would have
added a hundred times the machinery for nothing a visitor could perceive.

Change five things. Get the photograph first, because a personal brand page without
a photograph of the person is not a draft of the design, it is a draft of the idea.
Write a README on day one. Use a relative stylesheet path. Never write a colour
literally when a token exists, which would have made the rebrand a one-line change
instead of a six-line one with five missed. Write the real numbers into the HTML and
let the counter animate towards them.

Drop \file{netlify.toml}, or move its four security headers into the Caddy
configuration on the VPS where they would actually take effect, and drop the four
unused tokens.

\subsection{Open threads}

Eleven assets never delivered, the photograph above all; five blue literals on an
orange site; three dead stylesheet rules; seven stale comments; the canonical link
on the live ai.digitalise.agency pointing at this domain; and an undocumented
sibling relationship between two repositories that share a skeleton.

% ===========================================================================
\section{Glossary}
% ===========================================================================

\begin{description}[leftmargin=0pt,style=nextline]
\item[ARIA] Attributes beginning \code{aria-} describing a page to screen readers.
  They change nothing visually.
\item[Breakpoint] A screen width at which the layout changes. This project has two,
  at 960 and 600 pixels.
\item[Canonical URL] A declaration of where content authoritatively lives, used by
  search engines to avoid indexing it twice.
\item[Class] A label attached to elements so CSS and JavaScript can find them.
\item[clamp] A CSS function taking a minimum, a preferred value and a maximum, used
  so headings scale with the window between fixed limits.
\item[CSS] Cascading Style Sheets: the file controlling appearance.
\item[Custom property] A named value in CSS, written \code{-{}-name} and read with
  \code{var(-{}-name)}. The ``tokens'' are these.
\item[Data attribute] An attribute beginning \code{data-} that HTML ignores and
  JavaScript can read.
\item[Flexbox] A layout system arranging children in a row or a column.
\item[Grid] A layout system arranging children in rows and columns at once.
\item[HTML] HyperText Markup Language: content and structure.
\item[IntersectionObserver] A browser facility reporting when an element scrolls
  into view. All three scroll behaviours are built on it.
\item[JavaScript] The language running in the browser after the page loads.
\item[Media query] CSS applying only when a condition about the screen is true.
\item[object-fit] A CSS property controlling how an image fills a frame it does not
  exactly match. \file{ASSETS.md} specifies it for the missing photograph.
\item[order] A CSS property setting an item's position inside a flex or grid
  container, independent of its position in the HTML. It is what actually reverses
  the second case-study card.
\item[prefers-reduced-motion] A system setting asking pages to avoid animation,
  usually because it causes nausea or migraine. The stylesheet honours it; the
  JavaScript does not.
\item[Pseudo-element] An extra box CSS creates with no tag in the HTML, written
  \code{::before} or \code{::after}, used here for the hero's fading edges and the
  glow above the booking panel.
\item[rsync] A program copying files between computers, transferring only what
  differs.
\item[Selector] The part of a CSS rule saying which elements it applies to.
\item[SSH] Secure Shell: the standard way to connect to another computer securely.
\item[SVG] A picture described as shapes and coordinates rather than pixels, so it
  stays sharp at any size and can sit inside the HTML file.
\item[Token] In this project, one of the nineteen named values at the top of the
  stylesheet.
\item[VPS] Virtual Private Server: a rented computer controlled remotely. This
  project's is at \code{46.202.194.85}, which salmanadnan.com resolves to today.
\item[Workflow] A file describing a job GitHub runs automatically when something
  happens in the repository.
\end{description}

\vspace{1em}
\noindent
\textbf{For the full treatment}, including every line of every file with its
provenance, its connections and its failure modes, see
\file{docs/explainer/deep-dive.pdf} in this same folder.

\end{document}
